\begin{conclusion}

Ce mémoire étudie la reconnaissance de signes isolés de la Langue des Signes Malgache à partir de séquences vidéo. Il cherche à déterminer si des points de squelette extraits par MediaPipe permettent de comparer des modèles LSTM et GRU, puis d'intégrer le modèle retenu dans une application restituant le résultat sous forme de texte et de parole malgache.

Dans le cadre expérimental défini, le GRU atteint 90,0 \% d'exactitude et un F1-score macro de 0,897 sur les 100 séquences de test produites par un signeur absent de l'entraînement. Le LSTM atteint respectivement 87,0 \% et 0,864. Le GRU compte 137\,610 paramètres, contre 182\,538 pour le LSTM. Ces résultats motivent le choix du GRU pour le prototype. Ils concernent dix signes isolés et un seul signeur de test ; ils ne démontrent pas une reconnaissance générale de la LSM ni une traduction de phrases continues.

Le travail apporte un corpus de 400 séquences réparties entre quatre signeurs, un pipeline de préparation des données et une comparaison expérimentale de deux modèles récurrents. Il comprend également une application web composée d'un backend FastAPI utilisant le modèle converti au format TFLite et d'un frontend qui produit la synthèse vocale malgache dans le navigateur.

La validation de l'application confirme le fonctionnement des composants sur le scénario essayé, tout en mettant en évidence une erreur de classification sur la séquence de référence. L'effectif du corpus, le nombre de signeurs et la diversité limitée des conditions de capture restreignent la portée des conclusions. Le prototype doit donc être considéré comme une démonstration de faisabilité pour un vocabulaire fermé.

Les suites prioritaires sont d'élargir le corpus, d'évaluer le modèle sur plusieurs signeurs jamais vus et de tester l'application dans des conditions de capture variées. À plus long terme, des séquences annotées de signes continus permettraient d'étudier une traduction en contexte. La collaboration avec les signeurs malgaches sera nécessaire pour valider le vocabulaire et les résultats d'usage.

\end{conclusion}
